\documentclass[11pt]{article}

\usepackage[margin=1in]{geometry}
\usepackage{amsmath,amssymb,amsthm}
\usepackage{graphicx}
\usepackage{hyperref}
\usepackage{microtype}

\newtheorem{theorem}{Theorem}
\newtheorem{lemma}{Lemma}
\newtheorem{proposition}{Proposition}
\newtheorem{remark}{Remark}
\newcommand{\Fix}{\operatorname{Fix}}
\newcommand{\Lip}{\operatorname{Lip}}
\newcommand{\dist}{\operatorname{dist}}
\newcommand{\Id}{\operatorname{Id}}

\title{A Tube Construction for the Uniformly Lipschitz Unit-Ball Question in Decomposable Banach Spaces}
\author{Smoke-test packet for arXiv:2307.12958}
\date{June 12, 2026}

\begin{document}
\maketitle

\section*{Source Paper}

This packet concerns the paper
\[
\textit{Retraction methods and fixed point free maps with null minimal
displacements on unit balls}
\]
by Cleon S. Barroso and Valdir Ferreira, arXiv:2307.12958v2
(3 September 2024).  The stored source PDF is this arXiv v2, whose abstract
already uses the unit-ball/null-displacement framing.

The arXiv source paper asks the following questions after recalling the
Lin--Sternfeld positive minimal displacement theorem.

\begin{center}
\includegraphics[width=\textwidth]{figures/open_problem_crop.png}
\end{center}

The review thread also supplied the following screenshot of a
journal/published formulation.  This packet keeps the arXiv v2 PDF as the
source paper and includes the screenshot to document the exact stronger
unit-ball wording seen in review.

\begin{center}
\includegraphics[width=\textwidth]{figures/journal_question_crop.png}
\end{center}

\section*{Target}

We focus on the stronger unit-ball version suggested by the journal
screenshot:
\begin{quote}
For an infinite-dimensional Banach space \(X\), is there a uniformly
Lipschitz mapping \(T:B_X\to B_X\) satisfying
\(\mathrm d(T,B_X)=0\) and \(\Fix(T)=\varnothing\)?
\end{quote}

The result below answers this question for Banach spaces that split as a
topological direct sum of two infinite-dimensional closed subspaces.  This is
a partial result only; it does not settle the arbitrary Banach-space question.

\section*{Auxiliary Lipschitz Facts}

\begin{lemma}[Radial clipping]
For a normed space \(E\), the map
\[
\kappa_E(x)=
\begin{cases}
x, & \|x\|\le 1,\\
x/\|x\|, & \|x\|>1,
\end{cases}
\]
is \(2\)-Lipschitz from \(E\) onto \(B_E\).
\end{lemma}

\begin{proof}
Let \(\|x\|\le \|y\|\).  If both points are in \(B_E\), the assertion is
immediate.  If \(\|x\|\le1<\|y\|\), then
\[
\|x-y/\|y\|\|\le \|x-y\|+\|y-y/\|y\|\|
=\|x-y\|+\|y\|-1
\le 2\|x-y\|.
\]
If \(1<\|x\|\le\|y\|\), then
\[
\begin{aligned}
\left\|{\frac{x}{\|x\|}}-{\frac{y}{\|y\|}}\right\|
&\le \frac{1}{\|x\|}\|x-y\|
  +\|y\|\left|\frac{1}{\|x\|}-\frac{1}{\|y\|}\right|\\
&= \frac{1}{\|x\|}\|x-y\|
  +\frac{\|y\|-\|x\|}{\|x\|}
\le 2\|x-y\|.
\end{aligned}
\]
\end{proof}

\begin{lemma}[Variable-radius sphere selection]
Let \(Z\) be an infinite-dimensional Banach space, and let
\(R_Z:B_Z\to S_Z\) be a Lipschitz retraction.  Define, for \(r>0\),
\[
H(r,q)=r\,R_Z\bigl(\kappa_Z(q/r)\bigr),\qquad q\in Z.
\]
Then \(H\) is Lipschitz in the following uniform sense:
\[
\|H(r,q)-H(s,p)\|
\le 2\Lip(R_Z)\|q-p\|+(1+\Lip(R_Z))|r-s|
\]
for all \(p,q\in Z\) and \(r,s>0\).
\end{lemma}

\begin{proof}
For fixed \(r\),
\[
\|H(r,q)-H(r,p)\|
\le r\Lip(R_Z)\|\kappa_Z(q/r)-\kappa_Z(p/r)\|
\le 2\Lip(R_Z)\|q-p\|.
\]
It remains to compare the radii.  Assume \(0<r\le s\).  For fixed \(q\),
\[
\|H(r,q)-H(s,q)\|
\le |s-r|+r\Lip(R_Z)\|\kappa_Z(q/r)-\kappa_Z(q/s)\|.
\]
Let \(a=\|q\|\).  Since \(q/r\) and \(q/s\) lie on the same ray, a direct
case distinction gives
\[
r\|\kappa_Z(q/r)-\kappa_Z(q/s)\|\le s-r.
\]
Indeed the expression is \(0\) if \(a\ge s\), it is
\(a(s-r)/s\le s-r\) if \(a\le r\), and it is
\(r(1-a/s)\le r(1-r/s)\le s-r\) if \(r<a<s\).  This proves the stated
estimate.
\end{proof}

\begin{lemma}[Positive Lipschitz gauge with zero infimum]
If \(Y\) is infinite-dimensional, then there is a Lipschitz function
\(\eta:S_Y\to(0,1]\) such that \(\inf_{S_Y}\eta=0\).
\end{lemma}

\begin{proof}
By Riesz' lemma choose a sequence \((u_n)\subset S_Y\) and a number
\(\delta>0\) such that \(\|u_n-u_m\|\ge\delta\) whenever \(n\ne m\).  Let
\(\varepsilon_n\downarrow0\) with \(\varepsilon_n>0\), and fix
\(\lambda>0\).  Set
\[
\eta(u)=\min\left(1,\inf_n\{\varepsilon_n+\lambda\|u-u_n\|\}\right).
\]
This function is \(\lambda\)-Lipschitz and satisfies
\(\eta(u_n)\le\varepsilon_n\), hence its infimum is zero.  It is positive at
each point \(u\): at most one \(u_n\) can lie within \(\delta/2\) of \(u\),
while all the remaining terms are at least \(\lambda\delta/2\), and the
exceptional term, if present, is still strictly positive.
\end{proof}

\section*{Main Partial Result}

\begin{theorem}
Let \(X\) be an infinite-dimensional Banach space.  Suppose that
\[
X=Y\oplus Z
\]
as a topological direct sum, where \(Y\) and \(Z\) are both
infinite-dimensional closed subspaces.  Then there is a mapping
\(T:B_X\to B_X\) such that
\[
\Fix(T)=\varnothing,\qquad
\mathrm d(T,B_X)=0,\qquad
\sup_{n\ge1}\Lip(T^n)<\infty .
\]
\end{theorem}

\begin{proof}
Let \(P:X\to Y\) and \(Q:X\to Z\) be the associated bounded projections.
Choose constants \(a,b>0\) with \(a+b\le1\).  For instance one may take
\(a=1/2\) and \(b=1/4\).

By the Benyamini--Sternfeld theorem, the unit sphere of every
infinite-dimensional Banach space is a Lipschitz retract of its unit ball.
Let
\[
R_Y:B_Y\to S_Y,\qquad R_Z:B_Z\to S_Z
\]
be Lipschitz retractions.  Define a Lipschitz retraction
\[
A:Y\to aS_Y
\]
by
\[
A(y)=a\,R_Y\bigl(\kappa_Y(y/a)\bigr).
\]
It fixes \(aS_Y\), because if \(y=au\) with \(u\in S_Y\), then
\(\kappa_Y(y/a)=u\) and \(R_Y(u)=u\).

Let \(\eta:S_Y\to(0,1]\) be the positive Lipschitz gauge from the preceding
lemma.  For \(y\in aS_Y\), write \(u=y/a\in S_Y\) and set
\[
r(y)=b\,\eta(u).
\]
Then \(0<r(y)\le b\), \(r\) is Lipschitz on \(aS_Y\), and
\(\inf_{aS_Y}r=0\).

For \(x\in B_X\), put
\[
y(x)=A(Px),\qquad \rho(x)=r(y(x)),
\]
and define
\[
\Phi(x)=H(\rho(x),Qx)
=\rho(x)R_Z\bigl(\kappa_Z(Qx/\rho(x))\bigr).
\]
Finally set
\[
T(x)=y(x)-\Phi(x).
\]
The estimates above show that \(T\) is Lipschitz: \(x\mapsto y(x)\) is
Lipschitz, \(x\mapsto\rho(x)\) is Lipschitz, and the variable-radius lemma
controls \((\rho(x),Qx)\mapsto\Phi(x)\).

Since \(\|\Phi(x)\|=\rho(x)\le b\) and \(\|y(x)\|=a\), we have
\[
\|T(x)\|\le a+b\le1.
\]
Thus \(T\) maps \(B_X\) into \(B_X\).

Let also
\[
S(x)=y(x)+\Phi(x).
\]
We claim that
\[
T^2=S,\qquad T^3=T.
\]
Indeed, write \(y=y(x)\), \(\rho=\rho(x)\), and
\(\Phi(x)=\rho z\) with \(z\in S_Z\).  Since \(P(y-\rho z)=y\) and
\(A(y)=y\), the base and radius are unchanged after applying \(T\).  Moreover
\[
Q(y-\rho z)/\rho=-z\in S_Z,
\]
so \(\kappa_Z(-z)=-z\) and \(R_Z(-z)=-z\).  Hence
\[
T(Tx)=y+\rho z=S(x).
\]
The same calculation applied to \(S(x)=y+\rho z\) gives \(T(Sx)=y-\rho z\).
Thus the iterates of \(T\) alternate between the two Lipschitz maps \(T\) and
\(S\) for all \(n\ge1\), and consequently
\[
\sup_{n\ge1}\Lip(T^n)<\infty .
\]

There is no fixed point.  If \(T(x)=x\), then the calculation above gives
\[
x=T(x)=y-\rho z
\quad\text{and}\quad
T^2(x)=y+\rho z.
\]
But \(T^2(x)=T(x)\) would follow from \(T(x)=x\), so \(2\rho z=0\), impossible
because \(\rho>0\) and \(z\in S_Z\).

Finally, choose the separated sequence \((u_n)\subset S_Y\) used in the
construction of \(\eta\), and fix \(z_0\in S_Z\).  Put
\[
x_n=a u_n+b\eta(u_n)z_0.
\]
Then \(x_n\in B_X\), and the defining formula gives
\[
T(x_n)=a u_n-b\eta(u_n)z_0.
\]
Therefore
\[
\|T(x_n)-x_n\|=2b\eta(u_n)\le2b\varepsilon_n\to0.
\]
This proves \(\mathrm d(T,B_X)=0\).
\end{proof}

\section*{Interpretation}

This gives a new partial answer to the journal-style Q2 for every Banach
space that decomposes into two complemented infinite-dimensional pieces.  The
proof does not use a basis, unconditionality, rotations, uniform convexity,
or Hilbertian structure.  It uses only:
\begin{itemize}
\item a complemented splitting \(X=Y\oplus Z\);
\item the Benyamini--Sternfeld Lipschitz retraction from a unit ball onto its
unit sphere in each infinite-dimensional summand;
\item a variable-radius tube whose fiber map is period two after the first
iterate.
\end{itemize}

\section*{Why This Is Still Partial}

The complemented splitting is essential in the proof.  It supplies stable
base and fiber coordinates: the base point \(y(x)\) must remain unchanged
after the first iterate, while the \(Z\)-fiber changes sign.  Without a
projection or a Lipschitz section of a quotient map, there is no evident way
to recover the base coordinate from a point in the tube.  Thus the argument
does not cover hereditarily indecomposable spaces or arbitrary Banach spaces
with no infinite-dimensional complemented decomposition.

\section*{Addendum: A Universal \(\eta\)-Approximation}

There is also a useful universal observation in the opposite direction.  It
does not solve Q2, because the map depends on \(\eta\), but it shows that the
minimal displacement obstruction can always be made arbitrarily small while
retaining uniformly Lipschitz iterates.

\begin{proposition}
Let \(X\) be an infinite-dimensional Banach space.  For every \(\eta>0\) there
is a fixed-point-free mapping \(T_\eta:B_X\to B_X\) such that
\[
\sup_{n\ge1}\Lip(T_\eta^n)<\infty
\qquad\text{and}\qquad
\mathrm d(T_\eta,B_X)<\eta .
\]
\end{proposition}

\begin{proof}
Choose \(e\in S_X\) and \(f\in S_{X^*}\) with \(f(e)=1\), and put
\(Y=\ker f\).  Let \(0<\rho<1\).  Define
\[
\Phi:\rho B_Y\to S_X,\qquad
\Phi(y)=\frac{e+y}{\|e+y\|}.
\]
Its image is a small spherical cap \(C_\rho\).  Since
\(1\le\|e+y\|\le1+\rho\) for \(y\in\rho B_Y\), the map \(\Phi\) is
bi-Lipschitz onto \(C_\rho\), with constants bounded for small \(\rho\).
Moreover \(\operatorname{diam}(C_\rho)\to0\) as \(\rho\downarrow0\).

By Benyamini--Sternfeld, choose a Lipschitz retraction
\[
R_Y:B_Y\to S_Y.
\]
On \(\rho B_Y\), define
\[
W(y)=-\rho\,R_Y(y/\rho).
\]
Then \(W(\rho B_Y)\subset \rho S_Y\).  Once \(y\in\rho S_Y\), the retraction
identity gives \(W(y)=-y\).  Hence the iterates of \(W\) alternate between the
two Lipschitz maps
\[
y\mapsto -\rho R_Y(y/\rho),
\qquad
y\mapsto \rho R_Y(y/\rho),
\]
so \(\sup_n\Lip(W^n)<\infty\).  Also \(W\) has no fixed point: if
\(W(y)=y\), then \(\|y\|=\rho\), and consequently \(y=-y\), impossible.

Conjugate \(W\) to the cap:
\[
V=\Phi W\Phi^{-1}:C_\rho\to C_\rho .
\]
Then \(V\) is fixed-point free and uniformly Lipschitz under iteration.

It remains to retract \(B_X\) onto \(C_\rho\).  Choose
\[
a=\frac{1}{2(1+\rho)}.
\]
For \(x\in B_X\), set
\[
z(x)=\frac{x-f(x)e}{\max\{f(x),a\}}\in Y,
\]
and let \(Q_\rho:Y\to\rho B_Y\) be radial clipping:
\[
Q_\rho(z)=
\begin{cases}
z, & \|z\|\le\rho,\\
\rho z/\|z\|, & \|z\|>\rho.
\end{cases}
\]
The denominator in \(z(x)\) is bounded below by \(a\), and radial clipping is
Lipschitz, so
\[
P(x)=\Phi(Q_\rho(z(x)))
\]
defines a Lipschitz map \(P:B_X\to C_\rho\).  If \(x=\Phi(y)\in C_\rho\), then
\[
f(x)=\frac{1}{\|e+y\|}\ge\frac{1}{1+\rho}>a,
\qquad
\frac{x-f(x)e}{f(x)}=y,
\]
so \(P(x)=x\).  Thus \(P\) is a Lipschitz retraction.

Finally set
\[
T_\eta=V P
\]
after choosing \(\rho\) so small that \(\operatorname{diam}(C_\rho)<\eta\).
Since \(P|_{C_\rho}=I\), we have
\[
T_\eta^n=V^nP\qquad(n\ge1),
\]
and therefore the iterates are uniformly Lipschitz.  If \(T_\eta(x)=x\), then
\(x\in C_\rho\) and \(V(x)=x\), a contradiction.  Lastly,
\[
\mathrm d(T_\eta,B_X)
\le \sup_{x\in C_\rho}\|Vx-x\|
\le \operatorname{diam}(C_\rho)
<\eta .
\]
\end{proof}

\begin{remark}
This addendum is not a replacement for Q2.  It proves that for every positive
threshold \(\eta\) one may choose a fixed-point-free uniformly Lipschitz
dynamical system whose minimal displacement is below \(\eta\).  Q2 asks for
one single map with minimal displacement equal to zero.  The missing step is
still a countable gluing or invariant-skeleton construction that places
arbitrarily small caps inside one uniformly controlled map.
\end{remark}

\section*{Verification Notes}

No computational verification was used.  The proof is analytic.  The main
human-review points are:
\begin{itemize}
\item the \(2\)-Lipschitz estimate for radial clipping;
\item the uniform variable-radius estimate for
  \(H(r,q)=rR_Z(\kappa_Z(q/r))\);
\item the iterate identity \(T^2=S\), \(T^3=T\);
\item the fixed-point-free argument using \(\rho(x)>0\);
\item the sequence \(a u_n+b\eta(u_n)z_0\) witnessing null minimal
  displacement.
\end{itemize}

\section*{References Used}

The transfer and construction use Benyamini--Sternfeld's theorem that the
unit sphere of every infinite-dimensional Banach space is a Lipschitz retract
of its unit ball:
\[
\text{Y. Benyamini and Y. Sternfeld, Proc. Amer. Math. Soc. 88 (1983),
389--391.}
\]

\end{document}
